% Préambule commun des schémas TikZ — charte Liouaï
% Chaque schéma commence par % Préambule commun des schémas TikZ — charte Liouaï
% Chaque schéma commence par % Préambule commun des schémas TikZ — charte Liouaï
% Chaque schéma commence par \input{_preambule} puis \begin{document}.
\documentclass[tikz,border=6pt]{standalone}
\usepackage[T1]{fontenc}
\usepackage[utf8]{inputenc}
\usepackage[scaled=0.95]{helvet}
\renewcommand{\familydefault}{\sfdefault}
\usepackage{amsmath,amssymb}
\usepackage[european]{circuitikz}
\usetikzlibrary{arrows.meta,positioning,calc,angles,quotes,decorations.pathreplacing,backgrounds}

% Couleurs du logo Liouaï Digital Education
\definecolor{navy}{HTML}{243743}
\definecolor{vert}{HTML}{28B78D}
\definecolor{vertfonce}{HTML}{167A5C}
\definecolor{vertclair}{HTML}{EAF7F2}
\definecolor{gris}{HTML}{A1A9AE}
\definecolor{gristexte}{HTML}{5C6870}
\definecolor{fond}{HTML}{F3F5F6}
\definecolor{jaune}{HTML}{E5C800}
\definecolor{jauneclair}{HTML}{FFFBD6}

\tikzset{
  every picture/.style={line cap=round, line join=round, color=navy, text=navy},
  fleche/.style={-{Stealth[length=2.6mm,width=2mm]}, line width=1.1pt},
  vecteur/.style={fleche, line width=1.4pt},
  titre/.style={font=\bfseries\small, text=navy},
  note/.style={font=\footnotesize, text=gristexte},
  boite/.style={draw=navy, line width=0.9pt, rounded corners=5pt, fill=fond, align=center, inner sep=6pt},
}
 puis \begin{document}.
\documentclass[tikz,border=6pt]{standalone}
\usepackage[T1]{fontenc}
\usepackage[utf8]{inputenc}
\usepackage[scaled=0.95]{helvet}
\renewcommand{\familydefault}{\sfdefault}
\usepackage{amsmath,amssymb}
\usepackage[european]{circuitikz}
\usetikzlibrary{arrows.meta,positioning,calc,angles,quotes,decorations.pathreplacing,backgrounds}

% Couleurs du logo Liouaï Digital Education
\definecolor{navy}{HTML}{243743}
\definecolor{vert}{HTML}{28B78D}
\definecolor{vertfonce}{HTML}{167A5C}
\definecolor{vertclair}{HTML}{EAF7F2}
\definecolor{gris}{HTML}{A1A9AE}
\definecolor{gristexte}{HTML}{5C6870}
\definecolor{fond}{HTML}{F3F5F6}
\definecolor{jaune}{HTML}{E5C800}
\definecolor{jauneclair}{HTML}{FFFBD6}

\tikzset{
  every picture/.style={line cap=round, line join=round, color=navy, text=navy},
  fleche/.style={-{Stealth[length=2.6mm,width=2mm]}, line width=1.1pt},
  vecteur/.style={fleche, line width=1.4pt},
  titre/.style={font=\bfseries\small, text=navy},
  note/.style={font=\footnotesize, text=gristexte},
  boite/.style={draw=navy, line width=0.9pt, rounded corners=5pt, fill=fond, align=center, inner sep=6pt},
}
 puis \begin{document}.
\documentclass[tikz,border=6pt]{standalone}
\usepackage[T1]{fontenc}
\usepackage[utf8]{inputenc}
\usepackage[scaled=0.95]{helvet}
\renewcommand{\familydefault}{\sfdefault}
\usepackage{amsmath,amssymb}
\usepackage[european]{circuitikz}
\usetikzlibrary{arrows.meta,positioning,calc,angles,quotes,decorations.pathreplacing,backgrounds}

% Couleurs du logo Liouaï Digital Education
\definecolor{navy}{HTML}{243743}
\definecolor{vert}{HTML}{28B78D}
\definecolor{vertfonce}{HTML}{167A5C}
\definecolor{vertclair}{HTML}{EAF7F2}
\definecolor{gris}{HTML}{A1A9AE}
\definecolor{gristexte}{HTML}{5C6870}
\definecolor{fond}{HTML}{F3F5F6}
\definecolor{jaune}{HTML}{E5C800}
\definecolor{jauneclair}{HTML}{FFFBD6}

\tikzset{
  every picture/.style={line cap=round, line join=round, color=navy, text=navy},
  fleche/.style={-{Stealth[length=2.6mm,width=2mm]}, line width=1.1pt},
  vecteur/.style={fleche, line width=1.4pt},
  titre/.style={font=\bfseries\small, text=navy},
  note/.style={font=\footnotesize, text=gristexte},
  boite/.style={draw=navy, line width=0.9pt, rounded corners=5pt, fill=fond, align=center, inner sep=6pt},
}

\begin{document}
\begin{tikzpicture}[scale=0.8]
% 1 unités
\begin{scope}
\node[titre] at (2,3.7) {1. des unités différentes};
\draw[gris, ->] (0,0) -- (4.4,0) node[note, below left] {masse (g)}; \draw[gris, ->] (0,0) -- (0,2.8) node[note, right] {taille (m)};
\foreach \x/\y in {0.4/1.2,0.9/1.4,1.3/1.3,1.8/1.6,2.2/1.5,2.7/1.7,3.1/1.6,3.6/1.8}{\fill[navy] (\x,\y) circle (2pt);}
\draw[vertfonce, line width=1.4pt, ->] (0.2,1.45) -- (4.0,1.6);
\node[note, align=center] at (2,-1.0) {l'axe 1 suit la variable aux\\ plus grands nombres : il faut réduire};
\end{scope}
% 2 non linéaire
\begin{scope}[xshift=6.2cm]
\node[titre] at (2,3.7) {2. une structure courbe};
\draw[gris, ->] (0,0) -- (4.4,0); \draw[gris, ->] (0,0) -- (0,2.8);
\foreach \t in {0,15,...,180}{\fill[navy] ({2+1.6*cos(\t)},{0.6+1.6*sin(\t)}) circle (2pt);}
\draw[vertfonce, line width=1.4pt, ->] (0.1,1.25) -- (4.0,1.25);
\node[note, align=center] at (2,-1.0) {un seul paramètre (l'angle) suffirait,\\ mais aucun axe droit ne le capte};
\end{scope}
% 3 petite variance
\begin{scope}[xshift=12.4cm]
\node[titre] at (2,3.7) {3. l'essentiel est sur l'axe 2};
\draw[gris, ->] (0,0) -- (4.4,0); \draw[gris, ->] (0,0) -- (0,2.8);
\foreach \x in {0.3,0.8,...,3.8}{\fill[navy] (\x,1.7) circle (2pt); \filldraw[fill=jaune, draw=navy] (\x+0.25,1.0) circle (2pt);}
\draw[vertfonce, line width=1.4pt, ->] (0.1,1.35) -- (4.1,1.35);
\node[note, align=center] at (2,-1.0) {deux groupes séparés selon la\\ direction de plus faible dispersion};
\end{scope}
\end{tikzpicture}
\end{document}
